\section{Recovery relations, graphs and Tanner codes}\label{sec:prelim}

\subsection{Recovery relations and the question}\label{sec:prelim-question}

We use the definitions of the fifth paper of the series~\cite[Section~2]{DirectedRecovery2026}. A \emph{recovery
relation} on a finite coordinate set $V$ assigns to every $q\in V$ a recovery set $N(q)\subseteq V$. A \emph{storage code}
over a finite alphabet $\mathcal A$ with $\abs{\mathcal A}\ge2$ is a set $\mathcal C\subseteq\mathcal A^V$ such that two
words of $\mathcal C$ that agree on $N(q)$ agree at $q$, for every $q$~\cite[Definition~2.1]{DirectedRecovery2026}; the
recovery functions may be nonlinear. Its \emph{rate} is $\abs{V}^{-1}\log_{\abs{\mathcal A}}\abs{\mathcal C}$.
The \emph{paired incidence graph} $E(N)$ is the bipartite graph on two copies $L_p$, $R_q$ ($p,q\in V$) of the
coordinates, with an edge $L_pR_q$ exactly when $p\in N(q)$~\cite[Definition~2.3]{DirectedRecovery2026}. The
\emph{paired distance} of $p$ is $\dist_{E(N)}(L_p,R_p)$, infinite if the two ports are disconnected; the pair
$L_p,R_p$ is not an extra edge ($L_pR_p$ is an edge exactly when $p\in N(p)$). Its two degrees are $\deg R_q=\abs{N(q)}$ and $\deg L_p=\abs{\{q:p\in N(q)\}}$. A sequence
of relations $N_j$ on $V_j$ with $\abs{V_j}\to\infty$ satisfies the \emph{paired-distance condition} if for every fixed
$r$ the number of coordinates of paired distance at most $r$ is $o(\abs{V_j})$.

\begin{question}[{\cite[Question~9.1]{DirectedRecovery2026}}]\label{prob:q91}
Fix a degree bound~$\Delta$. Let $N_j$ be recovery relations on finite coordinate sets $V_j$ with $\abs{V_j}\to\infty$.
Suppose their incidence graphs have maximum degree at most~$\Delta$ and satisfy the paired-distance condition. Must
every sequence of nonempty finite-alphabet storage codes $\mathcal C_j\subseteq\mathcal A_j^{V_j}$ satisfy
\[
  \limsup_{j\to\infty}\frac{\log_{\abs{\mathcal A_j}}\abs{\mathcal C_j}}{\abs{V_j}}\le\frac12?
\]
Here $\abs{\mathcal A_j}\ge2$; the alphabets may vary.
\end{question}

The fifth paper adds that ``the fixed binary alphabet is already an important open case''. The eleventh paper of the
series formalizes the question as \lean{MergeDepth.Q91At} (at a degree bound $\Delta$, all alphabets) and
\lean{MergeDepth.Q91BinaryAt} (binary codes)~\cite{MergeDepth2026}, on top of the fourth paper's Lean definition of
storage codes (\lean{MotionDepthCodes.IsStorageCode})~\cite{CurvedMotion2026} and the fifth paper's definitions of the
paired incidence graph (\lean{DirectedRecovery.pairedIncidenceGraph}) and of the paired-distance condition
(\lean{DirectedRecovery.HasPairedDistanceSeparation}). Our formal
results are negations of these statements.

Two classes of relations are known to obey the half bound, and they explain the shape of our construction. If the
incidence components separate the two ports of every coordinate, every storage code has at most
$\abs{\mathcal A}^{\lfloor\abs V/2\rfloor}$ words~\cite[Lemma~4.1]{DirectedRecovery2026}, a Lean-verified result there.
In particular this holds when the coordinates split into two classes and every recovery set of a coordinate lies in the
other class: the code is then determined by its restriction to either class. We call such relations \emph{perfectly
separated}. Our relations are perfectly separated locally, on balls of growing radius, but not globally: a global twist
can connect the two ports of a coordinate, though only through very long paths.

\subsection{Graphs, trees and quotients}\label{sec:prelim-graphs}

Graphs may be infinite but are always locally finite and without loops. For a graph $G$ we write $G^0$ for its vertex
set, $G^1$ for its edge set and $\st(w)\subseteq G^1$ for the set of edges at a vertex~$w$. Unless we say otherwise,
graphs are simple. The distance between two edges is their distance in the line graph: $\dist(e,e)=0$, and for $e\ne f$,
$\dist(e,f)=1+\min\{\dist_G(x,y):x\in e,\ y\in f\}$. It is a metric, and automorphisms of $G$ preserve it. For an edge
$e$ and $R\ge0$, the \emph{ball} $B(e,R)$ is the set of edges at distance at most $R$ from~$e$.

\begin{lemma}[Balls in a regular tree]\label{lem:ball}
\notformalized
Let $\Tree$ be the $(q+1)$-regular tree with $q\ge2$. Every ball of radius $R\ge1$ around an edge has
\[
  \beta_q(R)=\frac{2(q^{R+1}-1)}{q-1}-1
\]
edges, and $\beta_q(R)<q^{R+1}$ for $q\ge4$.
\end{lemma}
\begin{informalproof}
Let $e=\{u,v\}$. An edge $f\ne e$ has $\dist(e,f)\le R$ if and only if its endpoint nearer to $e$ is at distance at
most $R-1$ from $\{u,v\}$. Since $\Tree$ is a tree, $f$ is then the edge from that endpoint away from $e$, so the edges
of $B(e,R)$ other than $e$ correspond bijectively to the vertices at distance $1,\dots,R$ from $\{u,v\}$. There are $2q^i$
vertices at distance~$i$, so $\abs{B(e,R)}=1+2\sum_{i=1}^Rq^i=2(q^{R+1}-1)/(q-1)-1$. For $q\ge4$ this is less than
$\tfrac23q^{R+1}<q^{R+1}$.
\end{informalproof}

A group $\Gamma$ acting on $G$ by automorphisms acts \emph{without inversions} if no element exchanges the two ends of an
edge. If $\Gamma$ acts freely on $G^0$ and without inversions, it acts freely on $G^1$: an element fixing an edge fixes
or exchanges its ends. The \emph{displacement} of a group $\Gamma'$ acting on $G$ is
\[
  s(\Gamma')=\inf\{\dist(e,\gamma e):e\in G^1,\ \gamma\in\Gamma'\setminus\{1\}\}.
\]
For an automorphism $\gamma$ without fixed vertices put $\ell(\gamma)=\min_x\dist_G(x,\gamma x)\ge1$. If $\gamma e\ne e$,
$x\in e$ and $y\in\gamma e$, then $y$ is within distance $1$ of $\gamma x$, so $\dist_G(x,y)\ge\ell(\gamma)-1$ and
\begin{equation}\label{eq:disp}
  \dist(e,\gamma e)\ge1+\bigl(\ell(\gamma)-1\bigr)=\ell(\gamma).
\end{equation}
Hence the displacement of a group acting freely and without inversions is at least the least value of $\ell$ on its
nontrivial elements.

For a subgroup $\Gamma'\subseteq\Gamma$ acting freely and without inversions, the quotient $X=\Gamma'\backslash G$ has
the orbits of vertices and of edges as vertices and edges, and two edges of $G$ have the same image under the projection
$\pi\colon G\to X$ if and only if they differ by an element of~$\Gamma'$. If $s(\Gamma')\ge2$, then $X$ has no loops
and $\pi$ maps the edges at each vertex $w$ bijectively onto the edges at $\pi(w)$: two edges at $w$ are at distance
at most $1$, so they are not identified, and an edge with both ends in one orbit would be at distance at most $1$ from
its translate. If $s(\Gamma')\ge3$, then $X$ is simple: if $e=\{u,v\}$ and $\{u,\gamma v\}$ are edges with
$\gamma\in\Gamma'\setminus\{1\}$, then $\{u,\gamma v\}$ meets both $e$ and $\gamma e$, so $\dist(e,\gamma e)\le2$. So for $s(\Gamma')\ge2$ the projection is a covering map of graphs.
When $G=\Tree$ is a tree it is a universal covering: paths lift uniquely from any lift of their start, and a map from a
tree to $X$ that is a local isomorphism on stars lifts to any covering of $X$ once one vertex is lifted, by lifting the
unique paths from that vertex.

\subsection{Local codes, Tanner codes and local checks}\label{sec:prelim-tanner}

Let $\F$ be a field and $G$ a graph. A \emph{local datum} on $G$ assigns to each vertex $w$ a subspace
$C_w\subseteq\F^{\st(w)}$, the \emph{local code}. The \emph{Tanner code}~\cite{Tanner1981} is
\[
  T(G)=\{c\colon G^1\to\F \;:\; c|_{\st(w)}\in C_w\text{ for every }w\in G^0\};
\]
when $G$ is infinite its words may have infinite support. A \emph{local check} at $w$ is a function on $G^1$ supported in
$\st(w)$ whose restriction to $\st(w)$ lies in the dual code $C_w^\perp$ for the standard bilinear form. Let $D(G)$ be the
space of finite sums of local checks. For a finitely supported $h$ and any $c$ on $G^1$ put
$\langle h,c\rangle=\sum_eh(e)c(e)$.

\begin{lemma}[Checks annihilate Tanner words]\label{lem:valid}
\formal{AboveHalfRate/Tanner.lean}{AboveHalf.pairing_eq_zero_of_mem_iSup,AboveHalf.pairing_eq_zero_of_mem_downChecks}
For every $h\in D(G)$ and every $c\in T(G)$, $\langle h,c\rangle=0$.
\end{lemma}
\begin{outline}
For a local check $g$ at $w$, $\langle g,c\rangle$ is the pairing of $g|_{\st(w)}\in C_w^\perp$ with
$c|_{\st(w)}\in C_w$. The pairing is additive in~$h$, so induction over sums of local checks finishes the proof.
\end{outline}

The local checks at $w$ \emph{interpolate on $t$ coordinates} if for every set $S\subseteq\st(w)$ of at most $t$ edges
and every $f\colon S\to\F$ some local check at $w$ restricts to $f$ on~$S$.

\begin{lemma}[MDS local codes interpolate]\label{lem:mds}
\formal{AboveHalfRate/Tanner.lean}{AboveHalf.interpolates_of_minWeight}
If every nonzero word of $C_w$ has more than $t$ nonzero coordinates, the local checks at $w$ interpolate on $t$
coordinates.
\end{lemma}
\begin{outline}
Let $S\subseteq\st(w)$ with $\abs S\le t$ and $Q=\st(w)\setminus S$. Restriction $C_w\to\F^Q$ is injective, since a word
vanishing on $Q$ has at most $\abs S\le t$ nonzero coordinates. So the functional $c\mapsto-\sum_{e\in S}f(e)c(e)$ on
$C_w$ factors through it, by a left inverse of the restriction, and is represented on $\F^Q$ by a vector~$g'$. The
function equal to $f$ on $S$ and to $g'$ on $Q$ is a local check with the required values.
\end{outline}

\begin{lemma}[Reed--Solomon local codes]\label{lem:grs}
\notformalized
Let $\F=\F_q$, let $\st(w)$ have $q+1$ elements, identified with the points of $\Pone(\F_q)$, and let $0\le m\le q-1$.
Choose a nonzero vector $x_e\in\F_q^2$ on the line of each point $e$ and scalars $a_e\ne0$. The code
$\{(a_ef(x_e))_e : f\text{ a binary form of degree }m\}$ has dimension $m+1$, every nonzero word has at least $q+1-m$
nonzero coordinates, and its local checks interpolate on $t=q-m$ coordinates.
\end{lemma}
\begin{informalproof}
A nonzero binary form of degree $m$ over $\F_q$ vanishes at no more than $m$ points of $\Pone$, because each zero
contributes a distinct linear factor. So evaluation is injective, since $m<q+1$, and every nonzero word has at most $m$
zero coordinates. The last claim is \cref{lem:mds}.
\end{informalproof}

A group $\Gamma$ acting on $G$ acts on functions by $(\gamma c)(e)=c(\gamma^{-1}e)$. A local datum is
\emph{$\Gamma$-equivariant} if $\gamma C_w=C_{\gamma w}$ for all $\gamma$ and~$w$. Then $\gamma D(G)=D(G)$ and
$\gamma T(G)=T(G)$. For a subgroup $\Gamma'$ acting freely and without inversions, with finite quotient
$X=\Gamma'\backslash G$, the \emph{Tanner code of the quotient} is
\[
  T(X)=\{c\colon X^1\to\F \;:\; c\circ\pi\in T(G)\}.
\]
When the projection maps stars bijectively, $T(X)$ is the Tanner code of the local datum that the equivariant datum
induces on~$X$.
